\section{Conclusion}

Turkish-CBE provides a controlled framework for studying the relationship between logical correctness and
presentation stability in Turkish and English. Its contribution is the explicit connection between latent
ordering, Turkish linguistic realization, query geometry, name controls, and denominator-aware paired
analysis. The constructed collection contains 3,000 independent source-order identities before query and
presentation expansion; those identities, rather than derived prompt counts, define the clustering
structure.

The available construction and preprocessing results demonstrate the need to distinguish measured surface
zeros from lemma familiarity, verified tokenizer assets from silently substituted ones, and raw agreement
from correct consistency. The executed pilot evaluation shows these distinctions operating on real models: perfect raw
invariance produced by a constant responder, and a frontier model whose correct consistency collapses under
a single answer-preserving Turkish construction while its marginal English scores remain respectable. These
pilot observations are exploratory. Genuine frozen-test evaluation after human stimulus review remains
necessary to answer RQ1--RQ5 confirmatorily. Future reports should preserve the present definitions, expose
uncertainty and missingness, and interpret improvements as evidence about the tested coverage rather than
proof of competence over an entire problem family.
